\documentclass[dvipsnames, svgnames, x11names]{beamer}
\special{dvipdfmx:config z 0} %取消PDF压缩，加快速度，最终版本生成的时候最好把这句话注释掉
\usetheme{mytheme}
% \usepackage{float}
% \usepackage{algorithm,algorithmic}

\newtcolorbox{boxmine}[2][]{colbacktitle=red!10!white,
colback=blue!10!white,coltitle=red!70!black,
title={#2},fonttitle=\bfseries,#1}

\title{第五章}
\subtitle{计算思维与算法}
\author{张斌}
\institute{理工系}
\date{\the\year 年 \the\month 月}




\begin{document}  

\begin{frame}
  \titlepage
  \thispagestyle{empty}
  \addtocounter{framenumber}{-1}
\end{frame}

\begin{frame}{}
  \begin{center}
      {\Huge \cbf{目 \ 录}}
  \end{center}
  \tableofcontents[subsectionstyle=hide/hide/show]
  \thispagestyle{empty}
  \addtocounter{framenumber}{-1}
\end{frame}

\section{计算思维}

\begin{frame}{什么是计算思维？}
\begin{itemize}                                                   
  \item {\cbf{人类科学发现的三大支柱}}
    \begin{itemize}
      \item 理论科学
      \item 实验科学
      \item 计算科学
    \end{itemize}
  \item \cbf{人类认识世界和改造世界的三种思维}
    \begin{itemize}
      \item 逻辑思维
      
      以推理和演绎为特征，以数学学科为代表。
      
      \item 实证思维
      
      以观察和总结自然规律为特征，以物理等学科为代表。
      \item 计算思维
      
      以设计和构造为特征，以计算机学科为代表。
    \end{itemize}
\end{itemize}
\end{frame}

\begin{frame}
    \frametitle{}
    \begin{boxmine}{ここにTitleを書くのだ}
      吾輩は猫である。名前はまだない。\\
      どこで生れたか頓（とん）と見当がつかぬ。何でも薄暗いじめじめした所でニャーニャー泣いていた事だけは記憶している。吾輩はここで始めて人間というものを見た。しかもあとで聞くとそれは書生という人間中で一番獰悪（どうあく）な種族であったそうだ。
      \end{boxmine}
\end{frame}

\begin{frame}
    \frametitle{顺序结构}
    \begin{tbox}哈哈哈哈哈哈
    \end{tbox}

        
            \begin{example}
                哈哈哈哈哈哈
            \end{example}

            \begin{tcolorbox}[title=\textbf{Lemma},colback=SeaGreen!10!CornflowerBlue!10,colframe=RoyalPurple!55!Aquamarine!100!] 子数整数问题 Python代码
            \end{tcolorbox}

            
  \end{frame}

  \begin{frame}
      \frametitle{}
      \begin{block}{}
        哈哈哈哈哈哈
    \end{block}
  \end{frame}

  \begin{frame}[fragile]
    \frametitle{子数整数问题 Python代码}
\begin{codeblock}{子数整数问题}
  k = int(input("请输入k："))
  n = 10000
  while n <= 30000:
      sub1 = int(n / 100)
      sub2 = int(n / 10) % 1000
      sub3 = n % 1000
      if sub1 % k == 0 and sub2 % k == 0 and sub3 % k == 0:
          print(n)
          n = n + 1
      else:
          n = n + 1
\end{codeblock}
\end{frame}

\begin{frame}[fragile]
    \frametitle{}
    \begin{algorithm}[H]  
      \caption{Calculate $y = x/y$}
    \begin{algorithmic}[1]
      \Begin
      \State \textbf{input} $x,y$;
      \If {$y=0$}
      \State \textbf{print} Error;
      \Else
      \State $z \gets x/y$
      \State \textbf{return}($z$);
      \End
      \end{algorithmic}
      \end{algorithm}
\end{frame}


\end{document}
